\begin{trapbox}

	\begin{description}[style=nextline, leftmargin=2.8em, labelsep=0.5em, itemsep=0.35em]
		\item[\textbf{\textcolor{CTrap}{易错点一（忽略非负条件）：}}]
			直接对负数使用均值不等式，如 $x<0$ 时得 $x+\frac1x\ge 2$，实际上 $x+\frac1x\le -2$（最大值为 $-2$，应转化为正数再使用）。
		\item[\textbf{\textcolor{CTrap}{正确理解（非负前提）：}}]
			使用前必须检查 $a_i\ge 0$，若变量为负应通过变形转化为正数再使用。
		\item[\textbf{\textcolor{CTrap}{错因分析（前提疏忽）：}}]
			未牢记均值不等式的非负条件，误以为对所有实数成立。
	\end{description}

	\medskip
	\begin{description}[style=nextline, leftmargin=2.8em, labelsep=0.5em, itemsep=0.35em]
		\item[\textbf{\textcolor{CTrap}{易错点二（取等条件缺失）：}}]
			求最值仅放缩得到不等式，未验证等号能否成立，例如求 $x+\frac1x$（$x\ge 2$）的最小值，由均值得 $x+\frac1x\ge 2$，但等号需 $x=1$ 不可取，故最小值不是 $2$。
		\item[\textbf{\textcolor{CTrap}{正确理解（验证等号）：}}]
			应用均值不等式求最值后必须令等号条件成立，并确认变量在定义域内，若不可取等则最值需另求。
		\item[\textbf{\textcolor{CTrap}{错因分析（忽略取等）：}}]
			把不等式放缩当成最值，未检查等号对应变量是否存在。
	\end{description}

	\medskip
	\begin{description}[style=nextline, leftmargin=2.8em, labelsep=0.5em, itemsep=0.35em]
		\item[\textbf{\textcolor{CTrap}{易错点三（定值构造错误）：}}]
			对 $y= x + \frac{1}{x^2}$ 直接使用均值，$x+ \frac1{x^2} \ge 2\sqrt{\frac1x}$，右边不是常数，无法得出最值。
		\item[\textbf{\textcolor{CTrap}{正确理解（和定或积定）：}}]
			使用均值不等式前必须通过拆项、配凑使各项和为定值或积为定值。
		\item[\textbf{\textcolor{CTrap}{错因分析（条件不全）：}}]
			对均值不等式“定值”条件理解不透，未先构造定值。
	\end{description}

\end{trapbox}
